\documentclass[onecolumn]{hipatia}
\usepackage{blindtext}
\newcommand{\superau}{\textsuperscript{\underline{a}}~}
\title{Inteligência Artificial}
\subtitle{Epístola}
\author{}
\begin{document}
\setcounter{page}{\epistolapage}
\maketitle
\leftskip=2.4cm
\rightskip=2.4cm

% \vspace{0.4cm}
\noindent Caro Leitor,
\vspace{0.4cm}

\noindent Que efeitos a Inteligência Artificial terá sobre a Matemática e, mais amplamente, sobre a sociedade que estamos construindo? Essa pergunta, que até pouco tempo parecia assunto de ficção, hoje atravessa nossas salas de aula, nossos laboratórios, nossas profissões e até nossa maneira de ver e produzir imagens.

\noindent É justamente desse encontro entre técnica, ciência, imagem e vida social que trata o primeiro artigo desta edição, em nossa seção \textsc{História}. Em \textit{Ciências e Fotografia: 500 anos de História (II)}, Ana Lucia Pinheiro Lima conclui seu vasto painel sobre fotografia, industrialização, comunicação de massa, arte e ciência, até chegar às transformações contemporâneas mediadas por algoritmos matemáticos.
Na seção \textsc{Memória}, revisitamos a trajetória de Maria Augusta de Araújo Moreno, professora marcante na formação matemática baiana e personagem central no processo de modernização do ensino no Colégio de Aplicação e no então Instituto de Matemática e Física da UFBA.
Em \textsc{Panorama}, Yuri Lima nos convida a um passeio por um dos cenários mais fascinantes dos sistemas dinâmicos: os bilhares. Entre modelos simples e regimes caóticos, o texto apresenta ideias centrais da área e mostra como perguntas geométricas aparentemente concretas abrem caminho para teorias profundas.
Em \textsc{Técnica}, Rafael Filipe dos Santos conduz o leitor por um roteiro elegante sobre números primos: da prova clássica de Euclides e variações em progressões aritméticas a resultados sobre fatores primos de sequências e distribuição dos primos, com exercícios instigantes ao final.
Já em \textsc{Problema}, Yure Carneiro e Samuel Feitosa reúnem uma seleção variada de desafios, com destaque para \textit{O Jogo do EL}, transitando por temas de álgebra, geometria, análise e combinatória em questões que convidam tanto à criatividade quanto ao rigor.

\noindent E, para terminar com uma nota bem-humorada: apesar de todas as incertezas sobre a IA, uma coisa já é certa: ela facilita bastante a vida do Editor. Seja na revisão de artigos, seja na criação de figuras. Na capa desta edição, por exemplo, pedimos ao ChatGPT que produzisse uma pintura no estilo de Edward Hopper representando matemáticos jogando bilhar em uma mesa elíptica.

\vspace{0.6cm}

\hfill Salvador, 25 de julho de 2026.

\hfill O Editor
\end{document}
